\chapter{Background and Literature Review}
\label{ch:background}

\section{Reinforcement Learning}
Reinforcement learning provides a mathematical framework for an agent learning
optimal behaviour through interaction. Its foundation is the Markov Decision
Process (MDP): a set of \emph{states} describing the environment, \emph{actions}
the agent may take, \emph{transition dynamics} mapping state--action pairs to next
states, and a \emph{reward} signal. The agent seeks a \emph{policy} --- a mapping
from states to actions --- that maximizes the expected cumulative discounted
return. This is captured by the Bellman equation~\cite{sutton2018}, which expresses
the value of a state as the immediate reward plus the discounted value of its
successor, enabling algorithms to \emph{bootstrap} value estimates from one another.
Classical methods such as Q-learning iteratively estimate the optimal action-value
function after each interaction, addressing the \emph{credit assignment} problem of
attributing outcomes to past actions.

\section{Deep Reinforcement Learning}
Tabular methods store a value for every state--action pair and become infeasible in
large or continuous state spaces. Deep Reinforcement Learning (DRL) replaces the
table with a neural-network function approximator. The Deep Q-Network (DQN)
pioneered this with \emph{experience replay} (sampling past transitions to break
temporal correlation) and \emph{target networks} (stabilizing the bootstrap target).

Whereas DQN is value-based, \emph{policy-gradient} methods directly optimize the
policy by ascending the expected return. Proximal Policy Optimization
(PPO)~\cite{schulman2017ppo} has become the standard modern choice: it constrains
how far the policy may move in a single update through a clipped surrogate
objective, preventing the destructive updates that destabilized earlier methods
while remaining simple and sample-efficient. Generalized Advantage Estimation
(GAE)~\cite{schulman2016gae} complements PPO by trading off bias and variance in
the advantage estimate through exponentially weighted temporal-difference errors;
it is now standard in policy-gradient implementations.

\section{Multi-Agent Reinforcement Learning}

\subsection{Fundamental Challenges}
Extending RL to multiple agents introduces complications absent in the
single-agent case. \textbf{Non-stationarity} is the central one: each agent
perceives the others as part of its environment, but those components evolve as the
others learn, violating the stationary-transition (Markov) assumption and voiding
single-agent convergence guarantees. \textbf{Credit assignment} also becomes
harder: when a team shares a reward, attributing it to individual actions is
difficult, and the joint action space grows exponentially (e.g.\ four agents with
four actions yield $4^4 = 256$ combinations). \textbf{Scalability} therefore
demands decomposing the global problem into local subproblems each agent can solve
from its own observation.

\subsection{Algorithmic Approaches}
\textbf{Independent Q-Learning (IQL)}~\cite{tan1993iql} is the simplest approach:
each agent learns its own value function, treating the others as environment and
ignoring non-stationarity entirely. It is simple and scalable, often works in
practice when mutual influence is limited, but lacks convergence guarantees and
struggles with cooperative credit assignment. Its policy-gradient analogue,
\textbf{Independent PPO (IPPO)}, applies PPO independently per agent and is the
``fully independent learners'' baseline used throughout this work.

\textbf{Value-decomposition} methods address credit assignment by learning how
individual values combine into a team value. \textbf{QMIX}~\cite{rashid2018qmix}
represents the joint action-value as a \emph{monotonic} mixture of individual
values, guaranteeing that improving an individual's value improves the team's,
while permitting decentralized execution.

These advances rest on the \textbf{Centralized Training with Decentralized
Execution (CTDE)} paradigm. Under CTDE, training may use privileged global
information (all agents' observations and actions) to stabilize learning and tame
non-stationarity, but the deployed policies act only on local observations,
preserving scalability and requiring no inter-agent communication at run time. A
useful analogy is a sports team: in practice a coach observes all players and
teaches coordination; in the match, players act independently on their local view,
applying what they learned.

\textbf{MADDPG}~\cite{lowe2017maddpg} pioneered CTDE for actor--critic methods with
centralized critics that see all agents' observations and actions during training
and decentralized actors at execution. \textbf{MAPPO}~\cite{yu2021mappo} combines
this centralized-critic architecture with PPO's robust optimization: a shared
centralized value function consumes all agents' observations during training, each
agent keeps its own decentralized policy, PPO clipping prevents destructive
updates, and GAE stabilizes advantages. MAPPO is reported to match or beat more
specialized cooperative algorithms across standard benchmarks, which is precisely
why it is the natural ``sophisticated cooperative method'' to test here.

Table~\ref{tab:algos} summarizes these algorithms.

\begin{table}[H]
\centering
\caption{Comparison of representative MARL algorithms.}
\label{tab:algos}
\small
\begin{tabular}{@{}llll@{}}
\toprule
\textbf{Algorithm} & \textbf{Training} & \textbf{Execution} & \textbf{Key characteristics} \\
\midrule
IQL / IPPO & Independent       & Decentralized & Simple, scalable; treats others as \\
           &                   &               & environment; unstable under non-stationarity \\[0.3em]
QMIX       & Centralized mixer & Decentralized & Monotonic value factorisation; \\
           &                   &               & cooperative settings only \\[0.3em]
MADDPG     & Centralized critic& Decentralized & Actor--critic CTDE; continuous actions; \\
           &                   &               & less stable, limited scalability \\[0.3em]
MAPPO      & Centralized critic& Decentralized & PPO-based; stable; strong coordination \\
           &                   &               & and scalability \\
\bottomrule
\end{tabular}
\end{table}

\section{Social Dilemmas in Multi-Agent Systems}
Social dilemmas describe situations where individually rational behaviour yields a
collectively poor outcome. The \emph{Prisoner's Dilemma} is canonical: both agents
benefit from mutual cooperation, yet each maximizes its own payoff by defecting
regardless of the other, so rational self-interest drives both to a worse outcome.
The \emph{Tragedy of the Commons} extends this to shared resources: each user
benefits from exploitation while the cost is distributed, creating pressure toward
overexploitation that destroys the resource for all. \emph{Sequential social
dilemmas}~\cite{leibo2017ssd} embed such incentives in grid-world environments over
many timesteps; the ``Harvest'' commons, where agents collect slowly regrowing
apples and over-harvesting permanently depletes a region, is a standard example.
Social dilemmas sit at the \emph{opposite} end of the spectrum from traffic and are
discussed as future work (Chapter~\ref{ch:futurework}).

\section{MARL for Traffic Signal Control}
Traffic signal control is naturally modelled as MARL: each traffic light is an
agent that observes local traffic state, selects a signal phase, and is rewarded
for traffic efficiency. Because intersections interact through traffic flow, agents
must learn policies that are good for the whole network, not just locally.
Traditional fixed-time and actuated controllers require manual tuning and cope
poorly with non-recurring congestion; RL promises adaptive controllers that learn
from experience. Early single-intersection deep-RL controllers improved on
classical methods but missed network effects, where a locally optimal decision
creates downstream problems.

\textbf{CoLight}~\cite{wei2019colight} addressed this by giving each agent
information about adjacent intersections and sharing network parameters across
agents, inducing coordinated behaviour despite decentralized execution. A
particularly effective family is \emph{pressure-based} control, which minimizes the
imbalance between incoming and outgoing vehicle counts per movement and connects to
max-pressure control from transportation theory. A large-scale
comparison~\cite{chu2020largescale} found that sharing neighbour information
significantly improves network performance \emph{but only with careful reward
design}: naive individual rewards, where each intersection minimizes only its own
delay, can harm overall throughput, whereas shared rewards help but demand careful
algorithm design to scale.

A central observation motivates this dissertation: from a system designer's
viewpoint traffic is \emph{fundamentally cooperative}. Network effects align
incentives --- gridlock hurts everyone --- which makes traffic an excellent,
controlled environment for asking whether an explicit coordination mechanism is
actually \emph{necessary}, or whether independent learners can recover the same
behaviour on their own.
